\documentclass[11pt]{article}
\usepackage[margin=1in]{geometry}
\usepackage{amsmath,amssymb}
\usepackage[T1]{fontenc}
\usepackage{lmodern}
\usepackage[hidelinks]{hyperref}

\title{Literature Status: Reflexive Isometry Range Need Not Be 1-Complemented}
\author{fa\_banach\_001}
\date{2026-06-15}

\begin{document}
\maketitle

\paragraph{Status.}
\texttt{literature\_already\_answered}.  This note records a later paper
that answers the reflexive form of an isometry-range question in
Randrianantoanina's survey \emph{Norm One Projections in Banach Spaces}
(arXiv:math/0110171).  No new proof is claimed here.

\paragraph{Source question.}
In Section ``Relations with isometries'', Question \texttt{qisometry}, the
survey asks whether, for a Banach space \(X\) and an isometry
\(T:X\to X\), the range \(T(X)\) must be \(1\)-complemented in \(X\).
The survey recalls that the unrestricted Banach-space question is already
negative for \(C[0,1]\), but says that the reflexive case remained open.
Equivalently, it was not known there whether every isometry in a reflexive
Banach space is a Wold isometry.

\paragraph{Supporting answer.}
Pelczar-Barwacz, \emph{A reflexive Banach space with an uncomplemented
isometric subspace} (arXiv:2309.15261), Theorem 0.1, constructs a reflexive
Banach space \(X\) with a basis \((e_i)_{i=1}^{\infty}\) such that the map
\(e_i\mapsto e_{2i}\) extends to an isometry from \(X\) onto
\[
Y=\overline{\operatorname{span}}\{e_{2i}:i\in\mathbb N\},
\]
while \(Y\) is not complemented in \(X\).  Hence \(Y\) is not
\(1\)-complemented in \(X\).  Composing the isometry \(X\to Y\) with the
inclusion \(Y\subset X\) gives an isometry of the reflexive Banach space
\(X\) into itself whose range is not \(1\)-complemented.

\paragraph{Scope.}
This gives a negative answer to the reflexive version singled out as open in
the source survey.  The broader nonreflexive formulation was already known
negative in the survey through the \(C[0,1]\) example, so the durable value of
this packet is the reflexive-case literature update.

\paragraph{Provenance evidence.}
The supporting paper explicitly frames its construction as answering the
question about isometries with non-\(1\)-complemented range and cites
Randrianantoanina's survey as related background.  A later 2024 paper
(arXiv:2407.15112) also states that Pelczar-Barwacz's construction yields an
isometry on a reflexive Banach space whose range is not \(1\)-complemented.

\begin{thebibliography}{9}
\bibitem{Randrianantoanina}
B. Randrianantoanina,
\emph{Norm One Projections in Banach Spaces},
arXiv:math/0110171.

\bibitem{PelczarBarwacz}
A. Pelczar-Barwacz,
\emph{A reflexive Banach space with an uncomplemented isometric subspace},
arXiv:2309.15261.

\bibitem{DeySain}
S. Dey and D. Sain,
\emph{Banach space contractions, isometric dilations and Wold isometries},
arXiv:2407.15112.
\end{thebibliography}

\end{document}
